% GENERATED FILE. Edit canonical lessons, then run npm run generate:books.
% canonical-source-hash: fnv1a64:7f1793161083ff7b
% canonical-lessons: GE-C531-aufwaermen, GE-C531-mikrowelle, GE-C531-wasserkocher, GE-C531-abwaschen, GE-C531-abtrocknen

\chapter{Aufwärmen und abwaschen}
\label{ch:ge-aufwarmen-und-abwaschen}
\hlchaptermodality{531}

\begin{chapteropening}
\textbf{By the end of this chapter:} \emph{I can talk about jobs in the kitchen.}

Complete the last lesson of chapter 531: I can talk about jobs in the kitchen.
\end{chapteropening}

\section[aufwärmen]{aufwärmen — to warm up, to reheat}

\label{lesson:GE-C531-aufwaermen}



\begin{quote}
Before the new one: say the German for bathrobe, then the German for pyjamas.
\end{quote}

\subsection*{You'll want to know: aufwärmen}
\textbf{aufwärmen} — \textquotedblleft{}to warm up, to reheat\textquotedblright{}.

\begin{quote}
Die Suppe von gestern wärmen wir heute auf.
\end{quote}

\begin{quote}
\emph{We're reheating yesterday's soup today.}
\end{quote}

Separable. Also in sport: \emph{sich aufwärmen} — to warm up.

\begin{cousinweb}
\emph{aufwärmen} is built on \emph{auf} \textquotedblleft{}up\textquotedblright{} and \emph{wärmen} \textquotedblleft{}to warm\textquotedblright{}, from \emph{warm}, inherited from Old High German \emph{warm}, from Proto-Germanic \emph{warmaz} (EtymDB).
\end{cousinweb}

\begin{grammarlens}[title={What you've built}]
One more word for telling, explaining and describing at B1.
\end{grammarlens}

\subsection*{Your turn}
\begin{itemize}
\raggedright
  \item \emph{Say it:} \emph{aufwärmen}
  \item \emph{Say it:} \emph{aufwärmen}, once more
  \item \emph{Say it:} a sentence of your own with \emph{aufwärmen}
  \item \emph{Recall:} say the German for bathrobe, then the German for pyjamas, then say \emph{aufwärmen} again
\end{itemize}

\begin{tcolorbox}[breakable,colback=teal!4,colframe=teal!35!black,title={Before you move on}]
What does aufwärmen mean? (\textquotedblleft{}To warm up, to reheat\textquotedblright{}.) Say it once more.
\end{tcolorbox}
\section[die Mikrowelle]{die Mikrowelle — microwave}

\label{lesson:GE-C531-mikrowelle}



\begin{quote}
Before the new one: say the German for pyjamas, then the German for to warm up.
\end{quote}

\subsection*{You'll want to know: die Mikrowelle}
\textbf{die Mikrowelle} — \textquotedblleft{}microwave\textquotedblright{}.

\begin{quote}
Er macht die Milch in der Mikrowelle warm.
\end{quote}

\begin{quote}
\emph{He heats the milk in the microwave.}
\end{quote}

Plural: \emph{die Mikrowellen}.

\begin{cousinweb}
\emph{Mikrowelle} is built on \emph{mikro-}, from Greek \emph{mikrós} \textquotedblleft{}small\textquotedblright{}, and \emph{Welle} \textquotedblleft{}wave\textquotedblright{} (EtymDB lists \emph{Welle}); it is usually explained as a loan translation of English \emph{microwave}.
\end{cousinweb}

\begin{grammarlens}[title={What you've built}]
One more word for telling, explaining and describing at B1.
\end{grammarlens}

\subsection*{Your turn}
\begin{itemize}
\raggedright
  \item \emph{Say it:} \emph{die Mikrowelle}
  \item \emph{Say it:} \emph{die Mikrowelle}, once more
  \item \emph{Say it:} a sentence of your own with \emph{die Mikrowelle}
  \item \emph{Recall:} say the German for pyjamas, then the German for to warm up, then say \emph{die Mikrowelle} again
\end{itemize}

\begin{tcolorbox}[breakable,colback=teal!4,colframe=teal!35!black,title={Before you move on}]
What does die Mikrowelle mean? (\textquotedblleft{}Microwave\textquotedblright{}.) Say it once more.
\end{tcolorbox}
\section[der Wasserkocher]{der Wasserkocher — kettle}

\label{lesson:GE-C531-wasserkocher}



\begin{quote}
Before the new one: say the German for to warm up, then the German for microwave.
\end{quote}

\subsection*{You'll want to know: der Wasserkocher}
\textbf{der Wasserkocher} — \textquotedblleft{}kettle\textquotedblright{}.

\begin{quote}
Der Wasserkocher ist kaputt, deshalb kochen wir das Wasser im Topf.
\end{quote}

\begin{quote}
\emph{The kettle is broken, so we're boiling the water in a pan.}
\end{quote}

Plural: \emph{die Wasserkocher}.

\begin{cousinweb}
\emph{Wasserkocher} is built on \emph{Wasser} \textquotedblleft{}water\textquotedblright{}, inherited from Old High German \emph{wazzar}, and \emph{Kocher} \textquotedblleft{}cooker\textquotedblright{}, from \emph{kochen}, an early borrowing from Latin \emph{coquere} \textquotedblleft{}to cook\textquotedblright{} (EtymDB).
\end{cousinweb}

\begin{grammarlens}[title={What you've built}]
One more word for telling, explaining and describing at B1.
\end{grammarlens}

\subsection*{Your turn}
\begin{itemize}
\raggedright
  \item \emph{Say it:} \emph{der Wasserkocher}
  \item \emph{Say it:} \emph{der Wasserkocher}, once more
  \item \emph{Say it:} a sentence of your own with \emph{der Wasserkocher}
  \item \emph{Recall:} say the German for to warm up, then the German for microwave, then say \emph{der Wasserkocher} again
\end{itemize}

\begin{tcolorbox}[breakable,colback=teal!4,colframe=teal!35!black,title={Before you move on}]
What does der Wasserkocher mean? (\textquotedblleft{}Kettle\textquotedblright{}.) Say it once more.
\end{tcolorbox}
\section[abwaschen]{abwaschen — to wash up, to do the dishes}

\label{lesson:GE-C531-abwaschen}



\begin{quote}
Before the new one: say the German for microwave, then the German for kettle.
\end{quote}

\subsection*{You'll want to know: abwaschen}
\textbf{abwaschen} — \textquotedblleft{}to wash up, to do the dishes\textquotedblright{}.

\begin{quote}
Nach dem Abendessen wäscht mein Bruder ab.
\end{quote}

\begin{quote}
\emph{After dinner my brother does the washing-up.}
\end{quote}

Separable, irregular: \emph{wusch ab, hat abgewaschen}. Also \emph{spülen}.

\begin{cousinweb}
\emph{abwaschen} is built on \emph{ab} \textquotedblleft{}off\textquotedblright{} and \emph{waschen} \textquotedblleft{}to wash\textquotedblright{}, inherited from Old High German \emph{wascan}, from Proto-Germanic \emph{waskaną} (EtymDB).
\end{cousinweb}

\begin{grammarlens}[title={What you've built}]
One more word for telling, explaining and describing at B1.
\end{grammarlens}

\subsection*{Your turn}
\begin{itemize}
\raggedright
  \item \emph{Say it:} \emph{abwaschen}
  \item \emph{Say it:} \emph{abwaschen}, once more
  \item \emph{Say it:} a sentence of your own with \emph{abwaschen}
  \item \emph{Recall:} say the German for microwave, then the German for kettle, then say \emph{abwaschen} again
\end{itemize}

\begin{tcolorbox}[breakable,colback=teal!4,colframe=teal!35!black,title={Before you move on}]
What does abwaschen mean? (\textquotedblleft{}To wash up, to do the dishes\textquotedblright{}.) Say it once more.
\end{tcolorbox}
\section[abtrocknen]{abtrocknen — to dry (dishes, yourself)}

\label{lesson:GE-C531-abtrocknen}



\begin{quote}
Before the new one: say the German for kettle, then the German for to wash up.
\end{quote}

\subsection*{You'll want to know: abtrocknen}
\textbf{abtrocknen} — \textquotedblleft{}to dry (dishes, yourself)\textquotedblright{}.

\begin{quote}
Ich wasche ab, und du trocknest ab.
\end{quote}

\begin{quote}
\emph{I wash and you dry.}
\end{quote}

Separable. \emph{Sich abtrocknen} — to dry yourself.

\begin{cousinweb}
\emph{abtrocknen} is built on \emph{ab} \textquotedblleft{}off\textquotedblright{} and \emph{trocknen} \textquotedblleft{}to dry\textquotedblright{}, from \emph{trocken} \textquotedblleft{}dry\textquotedblright{}, inherited from Middle High German \emph{trucken} (EtymDB).
\end{cousinweb}

\begin{grammarlens}[title={What you've built}]
One more word for telling, explaining and describing at B1.
\end{grammarlens}

\subsection*{Your turn}
\begin{itemize}
\raggedright
  \item \emph{Say it:} \emph{abtrocknen}
  \item \emph{Say it:} \emph{abtrocknen}, once more
  \item \emph{Say it:} a sentence of your own with \emph{abtrocknen}
  \item \emph{Recall:} say the German for kettle, then the German for to wash up, then say \emph{abtrocknen} again
\end{itemize}

\begin{tcolorbox}[breakable,colback=teal!4,colframe=teal!35!black,title={Before you move on}]
What does abtrocknen mean? (\textquotedblleft{}To dry (dishes, yourself)\textquotedblright{}.) Say it once more.
\end{tcolorbox}
